\section{From revealed paths to a permutation kernel}\label{sec:conditional}

We first specify the information being conditioned on and prove that the unobserved randomness acts on the whole free assignment. The distinction between a positional bijection and the labels carried by that bijection will matter when preparation intervals are repeated.

Identify the positions with $V=\mathbb F_2^d$, ordered lexicographically. Permutations send labels to positions, and $(gh)(x)=g(h(x))$. Put
\[
 R(x_1,\ldots,x_d)=(x_2,\ldots,x_d,x_1).
\]
One physical shuffle is $RS$, where $S$ independently fixes or exchanges each pair $\{x,x+e_1\}$. If $Y_t$ is the physical permutation and $X_t=R^{-t}Y_t$, then
\[
 X_{t+1}=R^{-t}S_{t+1}R^tX_t.
\]
Thus $X_t$ uses independent fair matching switches in directions $1,\ldots,d$ cyclically. A layer means one such update; a sweep consists of $d$ consecutive layers containing all directions. At sweep boundaries $X_t=Y_t$. All interval lengths used below end at such boundaries. Uniform measure on the full group, and on every ordered-injection space, is invariant under each layer.

\begin{lemma}[Conditional slot bijections]\label{lem:slot-kernel}
Fix a deterministic starting layout and a set of outside labels. For any feasible specification $e$ of their paths through a finite interval, conditional on $e$ all switch values on edges incident to an outside card are fixed, and all other switch values are independent fair bits. Let $F_s$ be the set of free positions at boundary $s$, and let $m=|F_s|$. Write $\iota_s:[m]\to F_s$ for the increasing bijection. Every completion of the unexposed bits induces a bijection $T_e:F_0\to F_T$, hence a permutation
\[
 g_e=\iota_T^{-1}T_e\iota_0\in S_m.
\]
Its conditional law depends on the outside paths and positions, but not on the incoming assignment of free labels. On consecutive intervals the corresponding permutations multiply with later intervals on the left. Conditional on all the outside paths, their remaining switch arrays are independent, so their conditional laws compose by convolution.
\end{lemma}
\begin{proof}
At a given layer, a specified outside path records both the incoming and outgoing endpoint of its matching edge and thereby fixes that edge's coin. If two outside cards meet, feasibility makes their requirements agree. Conversely, prescribing those values forces the outside paths: induction on the layer determines their positions regardless of all other switch values. The path event is exactly a cylinder in the independent time-edge coin array. Conditioning therefore leaves every unprescribed time-edge coordinate independent and fair, including coordinates earlier than a future prescribed path segment.

At a layer there are three cases. An outside--outside edge has no free slot. An outside--free edge has one free input and one free output, and the prescribed coin maps the former to the latter. A free--free edge has two free inputs and outputs, and its unexposed coin chooses the identity or transposition. These maps are bijections between the free input and output sets. Their composition is $T_e$. Relabelling the free cards changes neither the cylinder event nor these maps. If an incoming label assignment is a bijection $a:[m]\to\mathcal L$ from the slot indices to the free-label set $\mathcal L$, the outgoing assignment is $a\circ g_e^{-1}$. Consequently a uniform $g_e$ gives a uniform assignment, whatever $a$ is.

At an intermediate boundary the same increasing map $\iota_s$ is used as the terminal index map of one interval and the initial index map of the next. These maps cancel in the product, giving the stated multiplication order. The unexposed time-edge coordinates of disjoint intervals are disjoint. The cylinder description proves their conditional independence even when all outside paths, including future ones, have been revealed. Each factor law is the one obtained by conditioning only on that interval's outside paths and starting outside layout.
\end{proof}

The lemma extends the elementary path-cylinder observation to a law on $S_m$. It does not assert that this law is already uniform. In particular, its one-card marginals need not determine it. The next two sections supply the contraction of the full law.

\begin{figure}[ht]
\centering
\begin{tikzpicture}[>=Stealth, node distance=18mm]
\node (a) {$[m]$};
\node[right=24mm of a] (f) {$F_0$};
\node[right=30mm of f] (g) {$F_T$};
\node[right=24mm of g] (b) {$[m]$};
\draw[->] (a)-- node[above] {$\iota_0$} (f);
\draw[->] (f)-- node[above] {$T_e$} (g);
\draw[->] (g)-- node[above] {$\iota_T^{-1}$} (b);
\draw[->] (a.south) to[out=-45,in=-135] node[below] {$g_e=\iota_T^{-1}T_e\iota_0$} (b.south);
\end{tikzpicture}
\caption{A fixed revealed path environment determines the free sets. Unexposed switches determine a random bijection between them. Increasing slot indices turn this bijection into a permutation on one fixed group, so consecutive kernels can be multiplied.}
\label{fig:slots}
\end{figure}

For one tracked free card, Lemma~\ref{lem:slot-kernel} gives a simpler linear evolution. Its conditional masses are averaged across a free--free edge and transported across a mixed edge. The uniform vector on the free slots follows exactly the same transport. A signed difference from uniform therefore has sum zero and nonincreasing Euclidean norm. This linear evolution is adapted to the progressively revealed paths, despite being a formula for the law conditional on their entire final specification: later prescribed coins do not reveal earlier unexposed ones.

We will also use uniqueness of an individual path through a single sweep. Just before direction $i$ is updated, the earlier coordinates equal the card's terminal coordinates and the later coordinates equal its initial coordinates. Thus the initial and terminal positions determine every intermediate position. This statement holds for any cyclic phase and either order of the coordinates.

Finally, the support bound in Theorem~\ref{thm:methods} needs no conditioning. There are at most $2^{tn/2}$ coin strings in $t$ physical shuffles. Their support has uniform mass at most $2^{tn/2}/n!$, proving the bound. In particular distance at most $1/4$ requires
\[
 t\ge\left\lceil\frac2n\log_2\frac{3n!}{4}\right\rceil
       =2d-O(1).
\]
For comparison with the lower constants in the separate constructions, the elementary inequality $n!\ge(n/2)^{n/2}$ also gives distance at least $1-2^{-n/2}$ whenever $t\le d-2$. Stirling's bound $\log(n!)\ge n\log n-O(n)$ gives distance tending to one uniformly for $t\le d$. In particular the weaker lower constants $d/2$, $d-1$ and $d$ used by those constructions follow from this same support calculation.

Changing the initial deck translates a permutation law, so its distance from uniform is unchanged. Mixing exists in each fixed dimension: one sweep has positive probability of being the identity or any specified cube-edge transposition, and those transpositions generate $S_n$. Identity sweeps pad finite products to a common length, giving a positive uniform minorization for a sufficiently long sweep block.
